% =============================================================================
% 3.5 Discussion
% =============================================================================
\section{Discussion}
\label{sec:ch3-discussion}

\subsection{Forces du système}

Quatre propriétés distinguent la plateforme des outils recensés à la
section~\ref{sec:ch1-solutions}.

\begin{enumerate}
  \item \btpterm{Capitalisation active.} Le patrimoine n'est plus seulement
        archivé~: il est interrogeable et génératif.
  \item \btpterm{Traçabilité.} Chaque génération est ancrée dans des
        extraits identifiés, eux-mêmes rattachés à une pièce, une page et
        une version.
  \item \btpterm{Fiabilité procédurale.} La revue humaine n'est pas une
        option de l'interface~: c'est une étape du graphe d'exécution.
  \item \btpterm{Évolutivité.} L'architecture en couches permet d'ajouter
        un graphe ou un n\oe ud agentique sans reconstruire le workflow.
\end{enumerate}

\subsection{Limites actuelles}

Ces forces ne masquent pas quatre limites, qu'il faut énoncer clairement.

\begin{enumerate}
  \item \btpterm{Dépendance au corpus.} La qualité des réponses est bornée
        par la qualité et la couverture des documents patrimoniaux. Un
        bureau d'études dont la mémoire est pauvre n'obtiendra pas, par
        magie, des références qu'il n'a pas.
  \item \btpterm{OCR imparfait.} Les scans de mauvaise qualité dégradent
        l'indexation, donc la recherche, donc la génération. Le goulot
        d'étranglement mesuré à la section~\ref{sec:ch3-perfs} le confirme.
  \item \btpterm{Asymétrie linguistique.} Les modèles restent moins
        performants en français qu'en anglais, a fortiori sur le français
        juridique des marchés publics camerounais.
  \item \btpterm{Spécificité de certains dossiers.} Des consultations très
        atypiques, hors du patrimoine, dépassent les capacités du système,
        qui doit alors s'abstenir --- ce qu'il fait, mais au prix d'une
        aide réduite.
\end{enumerate}

\subsection{Pistes d'amélioration}

\begin{enumerate}
  \item \btpterm{Adaptation de domaine.} Un affinage ou, à défaut, des
        invites et des glossaires spécialisés réduiraient l'asymétrie
        linguistique.
  \item \btpterm{Récupération agentique localisée.} Encapsuler une
        recherche multi-sauts dans un n\oe ud du workflow, sans rendre
        l'ensemble agentique.
  \item \btpterm{Graphe de connaissances.} Phase 5, une fois le RAG hybride
        consolidé.
  \item \btpterm{Hygiène du corpus.} Outils de détection des doublons, des
        attestations périmées et des scans illisibles, en amont de
        l'indexation.
  \item \btpterm{Multimodalité.} Intégration des plans, schémas et images,
        aujourd'hui hors périmètre.
\end{enumerate}

\subsection{Conclusion du chapitre}

La plateforme démontre la faisabilité d'un système de capitalisation et de
génération assistée de dossiers d'appels d'offres dans le secteur du BTP
camerounais. Les résultats obtenus montrent qu'une recherche hybride
améliore significativement la pertinence, qu'une génération contrainte
produit des sections traçables, et qu'un workflow avec revue humaine
porte la conformité à un niveau compatible avec le caractère engageant
d'une offre. Les limites identifiées dessinent un programme de travail,
non un obstacle dirimant.

\btpEndMark
